\documentclass[12pt]{exam}

\usepackage[margin=1in]{geometry}
\usepackage{setspace}

\setstretch{1.15}

\nopointsinmargin
\pointformat{}

% Uncomment ONE of these:
\noprintanswers
% \printanswers

\begin{document}

% \begin{center}
%   {\Large \textbf{Personal Finance Final Exam}}\\
%   \vspace{0.2cm}
%   Name:\underline{\hspace{2.5in}} \\
%   \vspace{0.2cm}
%   Time Allowed: 75 Minutes
% \end{center}
\noindent
\begin{tabular*}{\textwidth}{@{\extracolsep{\fill}}lr}
  {\Large \textbf{Personal Finance Final Exam}} & Name:\underline{\hspace{2in}} \\
  & Date:\underline{\hspace{2in}}
\end{tabular*}

\vspace{0.3cm}

\noindent
75 Minutes,100 points plus 10 bonus points.
Up to 5 pages of notes and calculator permitted.
Answer all multiple-choice questions by selecting the best answer.
Answer all short-response questions clearly, using correct spelling, punctuation, and grammar.
Show all work for financial math questions.

\vspace{0.3cm}

\hrule
\vspace{0.4cm}

%%%%%%%%%%%%%%%%%%%%%%%%%%%%%%%%%%%%%%%%%%%%%%%%%%%%%%
\section*{Part I: Multiple Choice}

\begin{questions}

  \question[2]
  The tendency to attach more worth to something just because you already own it is called
  \begin{choices}
    \choice Opportunity cost
    \CorrectChoice Endowment effect
    \choice Diversification
    \choice Inflation
  \end{choices}

  \vspace{0.3cm}

  \question[2]
  Gross pay refers to:
  \begin{choices}
    \choice Income after taxes
    \CorrectChoice Income before deductions
    \choice Hourly wage only
    \choice Total yearly salary after bonuses
  \end{choices}

  \vspace{0.3cm}

  \question[2]
  FDIC insurance protects:
  \begin{choices}
    \choice Stock investments
    \CorrectChoice Bank deposits up to legal limits
    \choice Credit card purchases
    \choice Retirement accounts
  \end{choices}

  \question[2]
  A stock represents:
  \begin{choices}
    \choice A loan to a company
    \CorrectChoice A share of ownership in a company
    \choice A government bond
    \choice A mutual fund
  \end{choices}

  \vspace{0.3cm}

  \question[2]
  Compound interest means:
  \begin{choices}
    \choice Interest is calculated only on the initial principal
    \CorrectChoice Interest is calculated on both the initial principal and accumulated interest
    \choice Interest is paid only at the end of the investment period
    \choice Interest rates are compounded annually
  \end{choices}

  \vspace{0.3cm}

  \question[2]
  Diversification helps investors by:
  \begin{choices}
    \choice Guaranteeing profits
    \CorrectChoice Reducing return volatility
    \choice Avoiding taxes
    \choice Eliminating inflation
  \end{choices}

  \vspace{0.3cm}

  \question[2]
  Why is starting retirement savings early important?
  \begin{choices}
    \choice Taxes are eliminated
    \choice You can always retire earlier
    \CorrectChoice Compound growth has more time to work
    \choice Your employer will cover everything
  \end{choices}

  \vspace{0.3cm}

  \question[2]
  APR stands for:
  \begin{choices}
    \choice Annual Payment Rate
    \choice Average Principal Return
    \CorrectChoice Annual Percentage Rate
    \choice Applied Payment Rule
  \end{choices}

  \question[2]
  The main purpose of insurance is to:
  \begin{choices}
    \choice Build wealth
    \CorrectChoice Transfer financial risk
    \choice Avoid taxes
    \choice Raise credit scores
  \end{choices}

  \vspace{0.3cm}

  \question[2]
  Why are credit cards potentially dangerous for many people?
  \begin{choices}
    \choice They cannot be used online
    \choice They are only available for those age 21 and older
    \CorrectChoice High interest deb can accumulate quickly
    \choice They replace debit cards
  \end{choices}

  \question[2]
  The main purpose of an emergency fund is:
  \begin{choices}
    \choice Vacation savings
    \choice Retirement planning
    \CorrectChoice To cover unexpected expenses
    \choice For investing in stocks
  \end{choices}

  \vspace{0.3cm}

  \question[2]
  Which action best protects consumers from identity theft?

  \begin{choices}
    \choice Using the same password everywhere
    \choice Sharing debit card information with friends
    \CorrectChoice Monitoring financial accounts regularly
    \choice Posting personal information online
  \end{choices}

  \vspace{0.3cm}

  \question[2]
  Which benefit is commonly offered by employers in addition to salary?

  \begin{choices}
    \choice APR
    \CorrectChoice Health insurance
    \choice Credit score increase
    \choice Sales tax refund
  \end{choices}

  \vspace{0.3cm}

  \question[2]
  Which worker would most likely earn overtime pay?

  \begin{choices}
    \choice A salaried executive employee
    \choice A self-employed investor
    \CorrectChoice An hourly employee working more than 40 hours in a week
    \choice A retiree collecting Social Security
  \end{choices}

  \vspace{0.3cm}

  \question[2]
  Which savings vehicle generally offers the highest liquidity?

  \begin{choices}
    \choice Certificate of deposit (CD)
    \CorrectChoice Checking account
    \choice 401(k) account
    \choice Stock portfolio
  \end{choices}

  \vspace{0.3cm}

  \question[2]
  Which situation is most likely to trigger an overdraft fee?

  \begin{choices}
    \choice Depositing a paycheck
    \CorrectChoice Spending more than your checking account balance
    \choice Opening a savings account
    \choice Using direct deposit
  \end{choices}

  \vspace{0.3cm}

  \question[2]
  Which investment is generally considered the most volatile?

  \begin{choices}
    \choice Savings account
    \choice U.S. Treasury bond
    \CorrectChoice Stock in a single company
    \choice Certificate of deposit
  \end{choices}

  \vspace{0.3cm}

  \question[2]
  Which type of insurance would help replace income if someone becomes unable to work because of injury?

  \begin{choices}
    \choice Auto insurance
    \CorrectChoice Disability insurance
    \choice Renters insurance
    \choice Collision insurance
  \end{choices}

  \vspace{0.3cm}

  \question[2]
  Which loan typically has the shortest repayment period?

  \begin{choices}
    \choice Mortgage loan
    \choice Student loan
    \choice Auto loan
    \CorrectChoice Payday loan
  \end{choices}

  \question[2]
  Which of the following is an example of a fixed expense?

  \begin{choices}
    \choice Buying concert tickets
    \CorrectChoice Monthly rent payment
    \choice Eating at restaurants
    \choice Purchasing new clothes
  \end{choices}

  \vspace{0.3cm}

  \question[2]
  What is the primary purpose of a credit report?

  \begin{choices}
    \choice To calculate income taxes
    \choice To track investment returns
    \CorrectChoice To show a person's borrowing and repayment history
    \choice To determine insurance premiums only
  \end{choices}

  \vspace{0.3cm}

  \question[2]
  Which investment account often provides tax advantages specifically for retirement savings?

  \begin{choices}
    \choice Checking account
    \choice Auto loan
    \CorrectChoice Roth IRA
    \choice Debit card account
  \end{choices}

  \vspace{0.3cm}

  \question[2]
  Which of the following is most likely to increase the cost of car insurance premiums?

  \begin{choices}
    \choice Maintaining a clean driving record
    \choice Choosing a higher deductible
    \CorrectChoice Having multiple speeding tickets
    \choice Driving fewer miles each year
  \end{choices}

  \vspace{0.3cm}

  \question[2]
  Why is it generally risky to co-sign a loan for someone else?

  \begin{choices}
    \choice It automatically lowers your taxes
    \choice It improves your credit score immediately
    \CorrectChoice You become legally responsible if the other person fails to pay
    \choice It eliminates interest charges
  \end{choices}

  %%%%%%%%%%%%%%%%%%%%%%%%%%%%%%%%%%%%%%%%%%%%%%%%%%%%%%
  \newpage
  \section*{Part II: Short Response (10 points each)}

  \question[10]
  Taylor earns \$18 per hour and works 25 hours per week.

  Explain:
  \begin{itemize}
    \item Why gross income differs from take-home pay
    \item Two deductions Taylor may see on a paycheck
  \end{itemize}

  \fillwithlines{2in}

  \begin{solution}
    Gross pay is income before deductions. Take-home pay is lower because taxes and other deductions are removed. Common deductions include federal tax, state tax, Social Security, and Medicare.
  \end{solution}

  \vspace{0.4cm}

  \question[10]
  Jordan has \$2,000 saved and is deciding whether to keep it in checking, savings, or stocks.

  Explain one advantage and one disadvantage of each.

  \fillwithlines{2.5in}

  \begin{solution}
    Checking offers liquidity but low interest. Savings earns interest but limited growth. Stocks offer growth potential but involve risk.
  \end{solution}

  \vspace{0.4cm}

  \question[10]
  A recent graduate buys a new car using a long loan term because the monthly payment looks low.

  Identify two risks and explain what should have been considered first.

  \fillwithlines{2.5in}

  \begin{solution}
    Risks include paying more interest over time and owing more than the car is worth. They should consider total loan cost and affordability.
  \end{solution}

  \question[10]
  A gym offers a membership for only \$20 per month, but the contract includes cancellation fees and automatic renewal.

  Explain:
  \begin{itemize}
    \item One reason the offer might be misleading
    \item One step a consumer should take before signing
  \end{itemize}

  \fillwithlines{2.5in}

  \begin{solution}
    The low monthly price may hide additional costs such as cancellation fees or automatic renewal terms. A consumer should read the contract carefully before signing.
  \end{solution}

  \vspace{0.4cm}

  \question[10]
  A student receives \$500 for graduation and wants to spend it immediately on expensive clothes.

  Explain:
  \begin{itemize}
    \item One behavioral bias that may influence this decision
    \item One smarter alternative use for the money
  \end{itemize}

  \fillwithlines{2.5in}

  \begin{solution}
    Impulse buying or present bias may influence the decision. Better alternatives include saving, investing, or paying necessary expenses.
  \end{solution}

  \vspace{0.4cm}

  \question[10]
  Two job offers are available:

  \begin{itemize}
    \item Job A: \$19/hour with no benefits
    \item Job B: \$17/hour with health insurance and retirement matching
  \end{itemize}

  Identify two factors someone should consider before choosing between these jobs.

  \fillwithlines{2.5in}

  \begin{solution}
    Students may discuss benefits, long-term earnings, healthcare costs, retirement contributions, scheduling, or advancement opportunities.
  \end{solution}

  \vspace{0.4cm}

  \question[10]
  A person chooses the cheapest possible health insurance plan without reading the details and later discovers the deductible is extremely high.

  Explain:
  \begin{itemize}
    \item What a deductible is
    \item Why the cheapest premium is not always the best choice
  \end{itemize}

  \fillwithlines{2.5in}

  \begin{solution}
    A deductible is the amount paid before insurance begins covering expenses. Lower premiums may come with high deductibles and greater out-of-pocket costs.
  \end{solution}

  \vspace{0.4cm}

  \question[10]
  A student takes out multiple loans for purchases they do not truly need.

  Explain two possible long-term consequences of excessive borrowing.

  \fillwithlines{2.5in}

  \begin{solution}
    Possible consequences include debt accumulation, damaged credit, missed payments, reduced financial flexibility, and higher interest costs.
  \end{solution}

  %%%%%%%%%%%%%%%%%%%%%%%%%%%%%%%%%%%%%%%%%%%%%%%%%%%%%%
  \section*{Part III: Financial Math}

  \question[10]
  You deposit \$5,000 into a fund.  It yields compound interest of 2\% the first year,
  3\% the second year,
  and 4\% the third year.  How much will you have at the end of three years?

  \vspace{1.5in}

  \question[10]
  You buy a one-year zero coupon \$1,000 bond for \$956.00.  when you redeem the bond
  at maturity for \$1,000, what was your annual interest rate earned over the year?

  %%%%%%%%%%%%%%%%%%%%%%%%%%%%%%%%%%%%%%%%%%%%%%%%%%%%%%
  \section*{Bonus Question}

  \bonusquestion[10]

  What is one financial lesson from this course that you think will be most useful in your life?

  \fillwithlines{2in}

\end{questions}

\end{document}
